\subsection{Convex functions over an open convex set}

PROPOSITION 21. --- Let f be a convex function, defined over the non-empty open convex set X in the topological vector space E. Then f is continuous if, and only if, it is bounded above when restricted to some non-empty open subset U of X.

The condition is obviously necessary, we prove that it is sufficient. Let \(x_{0} \in X\) be a point such that f is bounded above in a neighbourhood V of \(x_{0}\) ; we show firstly that \(f\) is continuous at \(x_0\) . By translation, we can restrict ourselves to the case when \(x_0 = 0\) and \(f(x_0) = 0\) ; moreover we can suppose that the neighbourhood \(V\) is balanced (I, p.~7, prop. 4). Suppose that \(f(x) \leqslant a\) in \(V\) ; for every \(\varepsilon, 0 < \varepsilon < 1\) , we observe that if \(x \in \varepsilon V\) , then \(x / \varepsilon \in V\) and \(-x / \varepsilon \in V\) . Applying inequality (1) of II, p.~16 to the points \(x / \varepsilon\) and 0 and to the number \(\lambda = \varepsilon\) , we see that \(f(x) \leqslant \varepsilon f(x / \varepsilon) \leqslant \varepsilon a\) ; applying it to points \(x\) and \(-x / \varepsilon\) and the number \(\lambda = 1 / (1 + \varepsilon)\) , gives \(f(x) \geqslant -\varepsilon f(-x / \varepsilon) \geqslant -\varepsilon a\) . Thus \(f(x)\) is arbitrarily small in \(\varepsilon V\) , if \(\varepsilon\) is sufficiently small, and \(f\) is continuous at \(x = 0\) .

Now let \(y\) be some point of \(X\) ; since \(X\) is open, there is a number \(\rho > 1\) such that \(z = \rho y\) belongs to \(X\) . Let \(g\) be the homothety \(x \mapsto \lambda x + (1 - \lambda)z\) of centre \(z\) and ratio \(\lambda = 1 - \frac{1}{\rho}\) , which transforms 0 into \(y\) ; for every point \(g(x) \in g(V)\) , we have from (1)

\[
f (g (x)) \leqslant \lambda f (x) + (1 - \lambda) f (z) \leqslant \lambda a + (1 - \lambda) f (z).
\]

Thus \(f\) is bounded above in a neighbourhood of \(y\) and hence, by the first part, is continuous at \(y\) . The proposition is proved.

COROLLARY. --- Every convex function \(f\) defined over an open convex set \(X\) in \(\mathbf{R}^n\) is continuous in \(X\) .

We can suppose that X is not empty. Then there exist, in X, \(n + 1\) affinely independent points \(a_{i}\) ( \(0 \leqslant i \leqslant n\) ) and the convex envelope of these points, S, contains the open non-empty set formed of the points \(\sum_{i=0}^{n} \lambda_{i} a_{i}\) with \(0 < \lambda_{i} < 1\) for all i and \(\sum_{i=0}^{n} \lambda_{i} = 1\) . By II, p.~17, prop. 20, f is bounded above in S and therefore is continuous.

In a topological vector space of infinite dimensions there exist, in general, linear non-continuous forms (II, p.~80, exerc. 25) and thus convex functions that are not continuous at any point.

